% ==========================================================================
%  Chapter 8 — Related Work and Comparative Analysis
% ==========================================================================
\chapter{Related Work}
\label{ch:related}

\epigraph{%
  Every isolation technology is a point in a multi-dimensional
  trade-off space: compatibility, performance, security, and
  operational complexity.  Understanding where each point lies is
  prerequisite to choosing wisely.%
}{}

User-Mode Linux does not exist in isolation.  The past two decades have
produced a rich ecosystem of virtualization, sandboxing, and kernel
abstraction technologies, each occupying a different point in the
design space.  This chapter provides a detailed comparison with nine
related systems, analyzing the architectural decisions that differentiate
them from \UML{} and the trade-offs those decisions entail.  We conclude
with a comprehensive comparison table (Section~\ref{sec:related-table})
that summarizes the landscape.

For each system, we describe the core architecture, identify the key
design decisions that differentiate it from \UML{}, and assess the
implications for compatibility, performance, security, and operational
deployment.


% ------------------------------------------------------------------
\section{gVisor}
\label{sec:related-gvisor}
% ------------------------------------------------------------------

\subsection{Architecture}

gVisor~\cite{gvisor} is Google's user-space kernel, open-sourced in
2018.  Its core component, the \emph{Sentry}, is a process written
in Go that reimplements approximately 200 Linux system calls.  When
a containerized application issues a system call, it is intercepted
and handled by the Sentry rather than by the host Linux kernel.  This
interposition provides a security boundary: even if the application
exploits a bug in the Sentry's syscall implementation, it cannot
directly compromise the host kernel.

gVisor has evolved through three platform backends~\cite{gvisor-platforms}:

\begin{description}[style=nextline,leftmargin=2em]
  \item[ptrace (2018).]
    The original platform.  The Sentry attaches to application
    processes via \texttt{ptrace} and intercepts syscalls.  This is
    architecturally identical to \UML{}'s original mechanism, and it
    suffers from the same performance limitations: each syscall
    requires multiple context switches between the tracer (Sentry)
    and the tracee (application).

  \item[KVM (2018).]
    The Sentry runs in a lightweight KVM virtual machine.
    Application code runs in guest ring 3; the Sentry runs in guest
    ring 0.  Syscall traps are handled as VM exits.  This is
    conceptually similar to \UML{}'s \KVMvTwo{} backend, but the
    ``kernel'' being run in ring 0 is the Go Sentry, not the real
    Linux kernel.

  \item[Systrap (2023).]
    The newest platform~\cite{gvisor-systrap}.  It uses a combination
    of shared memory, per-thread signal handlers, and a dedicated
    ``stub'' process to handle syscalls without the overhead of
    \texttt{ptrace} or \KVM{}.  The application's syscall instruction
    triggers a \texttt{SIGSYS} signal (installed via seccomp-BPF),
    which vectors to a handler that dispatches the syscall to the
    Sentry via shared memory.  This achieves near-native performance
    for many workloads.
\end{description}

\subsection{Key difference: reimplementation vs.\ the real kernel}

The fundamental architectural difference between gVisor and \UML{}
is that gVisor \emph{reimplements} the Linux syscall interface while
\UML{} \emph{is} the Linux kernel.  This distinction has profound
consequences:

\begin{description}[style=nextline,leftmargin=2em]
  \item[Compatibility.]
    gVisor's Sentry implements a subset of the Linux ABI.  System
    calls that are not implemented, or whose semantics differ subtly
    from the kernel's implementation, can cause application failures.
    Google maintains a compatibility matrix documenting which syscalls
    are supported and at what level of fidelity.  \UML{}, by contrast,
    supports every syscall that the Linux kernel supports, because it
    \emph{is} the Linux kernel.  There is no compatibility matrix
    because there is no reimplementation.

  \item[Bug fidelity.]
    A bug found in gVisor is a bug in the Go Sentry code.  It may or
    may not correspond to a real kernel bug.  A bug found in \UML{}
    is a bug in the Linux kernel source tree, directly applicable to
    every architecture.  This makes \UML{} suitable for kernel
    development and testing in a way that gVisor fundamentally is not.

  \item[Security posture.]
    gVisor's reimplementation is a \emph{feature} from a security
    perspective: by not running the real kernel, it avoids exposing
    the real kernel's attack surface to untrusted containers.  The
    Sentry presents a smaller, memory-safe (Go) attack surface.
    \UML{}, running the real kernel, inherits the real kernel's
    vulnerabilities.  However, those vulnerabilities are contained
    within a host-level process, so exploitation achieves
    compromise of a userspace process, not the host kernel.
\end{description}

\subsection{Performance comparison}

gVisor with the systrap platform achieves syscall latencies of
approximately 1--3\,$\mu$s for simple syscalls~\cite{gvisor-systrap}.
\UML{} with the \Seccomp{} backend achieves approximately 300\,ns;
with \KVMvTwo{}, approximately 100\,ns.  For I/O-intensive workloads,
gVisor's overhead is dominated by the cost of reimplementing filesystem
and networking operations in Go, whereas \UML{} uses the kernel's
native C implementations.  For compute-intensive workloads where
syscalls are infrequent, both systems approach native performance.

\begin{figure}[ht]
\centering
\begin{tikzpicture}[
    layer/.style={
      draw=UMLGrid,
      fill=#1,
      minimum width=5cm,
      minimum height=0.9cm,
      rounded corners=3pt,
      font=\small,
      align=center
    },
    label/.style={
      font=\footnotesize\itshape,
      color=UMLMuted,
      anchor=south
    },
    arrow/.style={
      -{Stealth[length=5pt]},
      thick,
      color=UMLMuted
    }
  ]

  % gVisor stack (left)
  \node[label] at (-3.5,4.6) {gVisor};
  \node[layer=KernelGold!8] (gapp) at (-3.5,4)  {Application};
  \node[layer=UMLTeal!8]    (gsen) at (-3.5,2.8) {Sentry (Go, $\sim$200 syscalls)};
  \node[layer=UMLBlue!8]    (gplat) at (-3.5,1.6) {Platform (ptrace/KVM/systrap)};
  \node[layer=UMLBlue!15]   (ghost) at (-3.5,0.4) {Host Kernel};

  \draw[arrow] (gapp.south) -- (gsen.north);
  \draw[arrow] (gsen.south) -- (gplat.north);
  \draw[arrow] (gplat.south) -- (ghost.north);

  % UML stack (right)
  \node[label] at (3.5,4.6) {UML};
  \node[layer=KernelGold!8] (uapp) at (3.5,4)  {Guest Application};
  \node[layer=UMLGreen!8]   (uker) at (3.5,2.8) {UML Kernel (real Linux)};
  \node[layer=UMLTeal!8]    (ubak) at (3.5,1.6) {Backend (seccomp/kvm-v2)};
  \node[layer=UMLBlue!15]   (uhost) at (3.5,0.4) {Host Kernel};

  \draw[arrow] (uapp.south) -- (uker.north);
  \draw[arrow] (uker.south) -- (ubak.north);
  \draw[arrow] (ubak.south) -- (uhost.north);

  % Separator
  \draw[dashed,color=UMLGrid] (0,-0.1) -- (0,5.0);

\end{tikzpicture}
\caption[gVisor vs.\ UML architecture]{%
  Architectural comparison of gVisor and \UML{}.  gVisor interposes a
  Go-language Sentry that reimplements Linux syscalls.  \UML{} runs the
  actual Linux kernel, using a backend for syscall interception.  The
  layers are analogous but the kernel layer is fundamentally different:
  reimplementation (gVisor) vs.\ the real kernel (\UML{}).
}
\label{fig:related-gvisor-uml}
\end{figure}


% ------------------------------------------------------------------
\section{Firecracker}
\label{sec:related-firecracker}
% ------------------------------------------------------------------

\subsection{Architecture}

Firecracker~\cite{agache2020,firecracker} is a virtual machine monitor
(VMM) developed by Amazon Web Services (AWS) and open-sourced in 2018.
Written in Rust, it uses the KVM API to create lightweight virtual
machines (``microVMs'') with a minimal device model.  Firecracker is
the execution engine behind AWS Lambda and AWS Fargate.

A Firecracker microVM includes:

\begin{itemize}[leftmargin=2em]
  \item A KVM-based virtual CPU with full hardware-assisted
        virtualization (Intel VT-x / AMD-V).
  \item A minimal device model: virtio-net, virtio-block, a serial
        console, and a minimal i8042 keyboard controller.  No PCI bus,
        no USB, no GPU emulation.
  \item A custom boot protocol that loads the guest kernel directly
        into memory, bypassing BIOS/UEFI boot.
\end{itemize}

Firecracker achieves cold-start times of under 125\,ms and memory
overhead of approximately 5\,MB per microVM.

\subsection{Key difference: full VM vs.\ process}

Firecracker runs a \emph{complete virtual machine} with its own kernel
instance, guest physical memory, and emulated hardware.  The guest kernel
runs in the VM's ring 0 (guest mode), with full hardware privilege within
the VM.  Communication between the guest and the host occurs through
virtio devices and VM exits.

\UML{} runs the kernel as a \emph{process}.  There is no emulated hardware,
no virtio devices (in the conventional sense), and no guest physical memory
separate from the host's process address space.  The \UML{} kernel
accesses host resources through system calls, not through device emulation.

\begin{description}[style=nextline,leftmargin=2em]
  \item[Isolation model.]
    Firecracker provides hardware-enforced isolation: the guest cannot
    access host memory outside its designated physical memory region
    because the CPU's EPT/NPT page tables prevent it.  \UML{} provides
    software-enforced isolation: the guest cannot access host memory
    because the host kernel's process isolation prevents it.  Both
    provide strong isolation, but Firecracker's hardware enforcement
    is a stronger guarantee against kernel vulnerabilities in the host.

  \item[Use case.]
    Firecracker is designed for multi-tenant cloud isolation: running
    thousands of untrusted customer workloads on shared hardware.
    \UML{} is designed for kernel development, debugging, testing,
    and research---environments where the guest is not necessarily
    untrusted, but where observability and fast iteration matter more
    than multi-tenant security.

  \item[Hardware requirements.]
    Firecracker requires \texttt{/dev/kvm} and hardware virtualization
    extensions.  \UML{} with the \Seccomp{} or \texttt{ptrace} backend
    requires neither.

  \item[Host tool access.]
    A Firecracker guest is an opaque VM from the host's perspective.
    Debugging the guest kernel requires GDB remote protocol over serial.
    A \UML{} guest is a host process: \texttt{gdb}, \texttt{strace},
    \texttt{perf}, and every other host tool work directly.
\end{description}

\subsection{Performance}

For I/O throughput, Firecracker benefits from virtio's efficient
batched I/O model and KVM's hardware-accelerated instruction
execution.  For syscall-heavy workloads, the VM exit cost
(approximately 1--2\,$\mu$s per exit on modern hardware) is
comparable to \UML{}'s \Seccomp{} backend.  \UML{}'s \KVMvTwo{}
backend, which also uses KVM but avoids the full VM device model,
can achieve lower per-syscall latency because the exit handler is a
simple function call within the \UML{} process, not a return to
a separate VMM process.


% ------------------------------------------------------------------
\section{Linux Kernel Library (LKL)}
\label{sec:related-lkl}
% ------------------------------------------------------------------

\subsection{Architecture}

The Linux Kernel Library (\textsc{LKL})~\cite{purdila2010,lkl}
compiles the Linux kernel as a static library (\texttt{liblinux.a})
that can be linked into an ordinary application.  The application
calls kernel functions directly---\texttt{lkl\_sys\_open()},
\texttt{lkl\_sys\_read()}, etc.---and the kernel executes in the
same address space and at the same privilege level as the application.

\textsc{LKL} achieves this by:

\begin{enumerate}[leftmargin=2em]
  \item Implementing a minimal \texttt{arch/lkl/} port that maps
        kernel threading to host pthreads, kernel memory allocation
        to host \syscall{malloc}, and kernel timers to host timer
        APIs.
  \item Disabling the MMU.  There is no virtual memory translation;
        kernel and application share a flat address space.
  \item Disabling process scheduling.  There is no userspace; the
        kernel exists solely to provide its subsystems (VFS, networking,
        block layer) as library functions.
\end{enumerate}

\subsection{Performance}

\textsc{LKL}'s in-process execution model eliminates all
interposition overhead.  A ``syscall'' is a function call: the
application pushes arguments onto the stack, calls the kernel
function, and gets the result back immediately.  There is no context
switch, no signal delivery, no trap handler.  This makes \textsc{LKL}
extraordinarily fast for workloads that consist primarily of kernel
operations.

Fuzzing studies have reported approximately 60$\times$ throughput
improvement when using \textsc{LKL} compared to QEMU/KVM targets.
This number is the motivation for many of the v2 redesign's
performance optimizations.

\subsection{Limitations}

\textsc{LKL}'s speed comes at a fundamental cost:

\begin{description}[style=nextline,leftmargin=2em]
  \item[No user/kernel separation.]
    The application and the kernel share an address space.  A buffer
    overflow in the kernel can corrupt application data and vice versa.
    There is no protection boundary, so \textsc{LKL} cannot be used to
    study real syscall behavior where \syscall{copy\_from\_user} /
    \syscall{copy\_to\_user} semantics matter.

  \item[No process model.]
    \textsc{LKL} does not support multiple processes, scheduling, or
    signals.  Bugs that depend on process interaction, scheduling
    decisions, or signal delivery are invisible.

  \item[No SMP.]
    \textsc{LKL} is single-threaded from the kernel's perspective.
    Concurrency bugs, which are among the most dangerous kernel bugs,
    cannot be found or reproduced.

  \item[No realistic attack surface.]
    Because there is no user/kernel boundary, the attack surface that
    \textsc{LKL} exposes to a fuzzer is a strict subset of the real
    kernel's attack surface.  Entire classes of vulnerabilities---those
    involving user-pointer validation, memory mapping semantics, and
    inter-process communication---are absent.
\end{description}

\subsection{The LKL--UML unification effort}

Hajime Tazaki proposed unifying \textsc{LKL} into \UML{} in a series
of RFC patches~\cite{tazaki2019} that reached v8 before stalling.
The core idea was to make \textsc{LKL} a ``backend'' or operating
mode of \UML{}, sharing the same \texttt{arch/um/} code but with a
compilation option that strips away process isolation for maximum
speed.  This would allow a single codebase to serve both the
high-fidelity (\UML{}) and high-speed (\textsc{LKL}) use cases.

The v2 redesign's backend-ops abstraction lays the groundwork for
this unification.  A hypothetical \texttt{lkl} backend could implement
the \texttt{struct um\_backend\_ops} interface with trivial stubs
(no-op for process creation, direct function call for syscall
dispatch), providing \textsc{LKL}-like performance within the \UML{}
framework.

\begin{figure}[ht]
\centering
\begin{tikzpicture}[
    box/.style={
      draw=UMLGrid,
      fill=#1,
      minimum width=3.5cm,
      minimum height=2.5cm,
      rounded corners=4pt,
      font=\small,
      align=center
    },
    label/.style={
      font=\small\bfseries,
      color=UMLInk
    }
  ]

  % LKL
  \node[box=UMLRed!6] (lkl) at (0,0) {
    \textbf{Single address space}\\[3pt]
    No MMU\\
    No processes\\
    No scheduling\\[3pt]
    \textit{60$\times$ speed}
  };
  \node[label,above=0.2cm of lkl] {LKL};

  % UML
  \node[box=UMLGreen!8] (uml) at (5,0) {
    \textbf{Full kernel}\\[3pt]
    Process isolation\\
    SMP scheduling\\
    Real syscall path\\[3pt]
    \textit{Full fidelity}
  };
  \node[label,above=0.2cm of uml] {UML};

  % QEMU
  \node[box=UMLBlue!8] (qemu) at (10,0) {
    \textbf{Full VM}\\[3pt]
    Hardware emulation\\
    Device model\\
    Guest physical memory\\[3pt]
    \textit{Any architecture}
  };
  \node[label,above=0.2cm of qemu] {QEMU/KVM};

  % Spectrum arrow
  \draw[-{Stealth[length=6pt]},very thick,color=UMLMuted]
    (-2,-2.2) -- (12,-2.2);
  \node[font=\footnotesize,color=UMLMuted,anchor=north] at (-2,-2.3)
    {Speed};
  \node[font=\footnotesize,color=UMLMuted,anchor=north] at (12,-2.3)
    {Fidelity};
  \node[font=\footnotesize\itshape,color=UMLAmber,anchor=north] at (5,-2.3)
    {UML v2 target: both};

\end{tikzpicture}
\caption[Speed--fidelity spectrum]{%
  The speed--fidelity spectrum of kernel execution environments.
  \textsc{LKL} maximizes speed by eliminating all isolation.
  QEMU/KVM maximizes fidelity by emulating complete hardware.
  \UML{} v2 aims to approach \textsc{LKL}'s speed while
  retaining full kernel fidelity.
}
\label{fig:related-spectrum}
\end{figure}


% ------------------------------------------------------------------
\section{QEMU and KVM}
\label{sec:related-qemu-kvm}
% ------------------------------------------------------------------

\subsection{Architecture}

QEMU~\cite{bellard2005} is a full-system emulator that translates
guest instructions to host instructions via a technique called
Tiny Code Generator (TCG) dynamic binary translation.  When combined
with KVM~\cite{kivity2007}, QEMU uses hardware virtualization
extensions to execute guest code directly on the CPU, with the TCG
path used only for emulating devices and handling privileged
instructions.

QEMU/KVM is the standard tool for kernel development and testing.
It provides:

\begin{itemize}[leftmargin=2em]
  \item Full hardware emulation: PCI bus, USB, network interfaces,
        block devices, GPU (via virtio-gpu or QXL), serial ports,
        RTC, APIC, IOAPIC, and more.
  \item Multi-architecture support: x86, ARM, RISC-V, PowerPC,
        MIPS, s390x, and others.
  \item Snapshot support via QEMU savevm.
  \item GDB debugging via the GDB remote protocol.
  \item Live migration between hosts.
\end{itemize}

\subsection{Comparison with UML}

\begin{description}[style=nextline,leftmargin=2em]
  \item[Architecture support.]
    QEMU can emulate any supported architecture on any host.  A
    developer can test an ARM kernel on an x86 workstation.  \UML{}
    runs only the host's native architecture (x86-64 on an x86-64
    host).  For cross-architecture testing, QEMU is the only option.

  \item[Device emulation.]
    QEMU provides realistic device emulation, essential for testing
    device drivers.  \UML{}'s devices are virtual abstractions
    (\texttt{ubd} for block, \texttt{tuntap}/\texttt{vector} for
    networking) that do not correspond to real hardware.  For driver
    development, QEMU is necessary.

  \item[Startup time.]
    QEMU/KVM requires 2--5 seconds to boot a kernel to userspace.
    \UML{} boots in under 500\,ms.  For workflows that require many
    boot cycles (bisection, fuzzing, CI), this difference is
    significant.

  \item[Debugging experience.]
    QEMU exposes a GDB remote protocol server.  The developer connects
    \texttt{gdb} to this server and debugs the guest kernel over a
    socket.  This works but imposes limitations: conditional breakpoints
    are evaluated on the debugger side (requiring round-trips for each
    check), variable inspection requires translating through guest page
    tables, and the full \texttt{gdb} Python scripting API is
    available but slow.

    \UML{} provides native \texttt{gdb} access.  The kernel is a
    local process; \texttt{gdb} reads its memory directly.  Conditional
    breakpoints evaluate at native speed.  The debugging experience
    is qualitatively superior.

  \item[Host tool access.]
    A QEMU guest is isolated from the host.  Sharing files requires
    virtio-9p, NFS, or SSH.  \texttt{strace}, \texttt{perf}, and
    other host tools cannot observe the guest kernel.  A \UML{} guest
    is a host process: all host tools work directly, and
    \texttt{hostfs} provides zero-overhead filesystem sharing.

  \item[Hardware requirements.]
    QEMU with TCG (software emulation) requires no hardware
    virtualization, but is 10--100$\times$ slower than native execution.
    QEMU/KVM requires \texttt{/dev/kvm}.  \UML{} with \Seccomp{}
    runs at near-native speed without any hardware requirements.
\end{description}

\subsection{Complementary roles}

\UML{} and QEMU/KVM are not competitors; they serve different
purposes.  QEMU/KVM is essential for cross-architecture testing,
driver development, and production virtualization.  \UML{} is
superior for kernel debugging, rapid iteration, and environments
lacking \texttt{/dev/kvm}.  A comprehensive kernel development
workflow uses both.


% ------------------------------------------------------------------
\section{Containers}
\label{sec:related-containers}
% ------------------------------------------------------------------

\subsection{Architecture}

Linux containers (Docker, Kubernetes, LXC) use namespaces and
cgroups~\cite{soltesz2007,namespaces,cgroups-v2} to provide
isolated execution environments that share the host kernel.  Each
container has its own view of the PID namespace, network namespace,
mount namespace, and (optionally) user namespace, but all containers
execute system calls against the same kernel instance.

Containers are extraordinarily lightweight---startup times of tens
of milliseconds, memory overhead of kilobytes---because there is
no kernel to boot, no device model to initialize, and no memory
to dedicate to a separate kernel instance.

\subsection{The shared-kernel problem}

The shared kernel is both containers' greatest strength and their
greatest weakness:

\begin{description}[style=nextline,leftmargin=2em]
  \item[Strength.]
    Because all containers share the kernel, there is no duplication
    of kernel memory, no boot overhead, and no performance penalty for
    syscall dispatch.  A containerized application's \syscall{read}
    call goes directly to the kernel with zero interposition overhead.

  \item[Weakness.]
    A kernel vulnerability is a vulnerability in \emph{every}
    container on the host.  Container escape exploits
    (CVE-2019-5736: runc; CVE-2020-15257: containerd; CVE-2022-0185:
    filesystem) demonstrate repeatedly that namespace and cgroup
    isolation do not constitute a security boundary.  The kernel
    developers explicitly document that namespaces are not a
    security feature for untrusted code.
\end{description}

\subsection{Comparison with UML}

\UML{} provides per-instance kernel isolation.  Each \UML{} instance
runs its own kernel, so a kernel vulnerability in one instance cannot
affect another instance or the host.  This comes at a cost:

\begin{itemize}[leftmargin=2em]
  \item \UML{} instances consume more memory (kernel per instance).
  \item \UML{} instances have longer startup times (kernel boot).
  \item Syscalls in \UML{} incur interposition overhead.
\end{itemize}

For multi-tenant isolation of untrusted workloads, the industry
has moved toward combining containers with lightweight VMs
(Firecracker, Kata Containers) or user-space kernels (gVisor).
\UML{} occupies a middle ground: lighter than a full VM, heavier
than a container, but with genuine kernel isolation and the unique
advantage of running the real kernel for development and testing
purposes.


% ------------------------------------------------------------------
\section{Unikernels}
\label{sec:related-unikernels}
% ------------------------------------------------------------------

\subsection{Architecture}

Unikernels~\cite{madhavapeddy2013,kuenzer2021} are specialized,
single-purpose machine images that combine an application with the
minimal set of kernel libraries needed to run it.  The application
and kernel run in a single address space at a single privilege level,
eliminating the user/kernel boundary.

Notable unikernel projects include:

\begin{description}[style=nextline,leftmargin=2em]
  \item[MirageOS~\cite{madhavapeddy2013}.]
    Written in OCaml.  Compiles an application and its OS dependencies
    into a single Xen or KVM bootable image.  Extremely small attack
    surface (no shell, no multi-user system, no unnecessary drivers).

  \item[Unikraft~\cite{kuenzer2021}.]
    A modular unikernel construction toolkit.  The developer selects
    which kernel components (scheduler, memory allocator, network
    stack) to include.  Achieves boot times under 1\,ms for minimal
    configurations and throughput exceeding that of Linux for
    specialized workloads.

  \item[IncludeOS.]
    A C++ unikernel for cloud applications.  The application
    \texttt{\#include}s the OS as a library.
\end{description}

\subsection{Comparison with UML}

\begin{description}[style=nextline,leftmargin=2em]
  \item[Boot time.]
    Unikernels boot in milliseconds (sometimes sub-millisecond for
    Unikraft), compared to \UML{}'s $\sim$500\,ms.  However,
    unikernels typically run on a hypervisor (Xen or KVM), so the
    total startup includes hypervisor overhead.

  \item[Compatibility.]
    Unikernels sacrifice compatibility for specialization.  A
    MirageOS unikernel cannot run arbitrary Linux binaries.  Unikraft
    has a Linux-compatible syscall layer, but it is incomplete.
    \UML{} runs unmodified Linux binaries because it is the real
    Linux kernel.

  \item[Debugging.]
    Unikernels are notoriously difficult to debug.  There is no
    shell to log into, no \texttt{gdb} to attach, and the
    single-address-space design means that a crash destroys the
    entire environment including any in-memory logs.  \UML{} provides
    full native debugging.

  \item[Security model.]
    Unikernels reduce the attack surface by eliminating unnecessary
    code, but the single-address-space design means that any
    vulnerability that achieves code execution has full control
    of the machine.  \UML{} maintains user/kernel separation,
    so a compromised guest process is still constrained by the
    \UML{} kernel's process isolation.

  \item[Use case overlap.]
    There is minimal overlap between unikernels and \UML{}.
    Unikernels are for production deployment of specialized services.
    \UML{} is for kernel development, testing, and research.  They
    occupy opposite corners of the design space.
\end{description}


% ------------------------------------------------------------------
\section{Dune}
\label{sec:related-dune}
% ------------------------------------------------------------------

\subsection{Architecture}

Dune~\cite{belay2012} is a system that provides safe user-level
access to privileged CPU features.  Developed at Stanford, Dune uses
Intel VT-x to place a process into a non-root mode VM where it has
access to ring 0 features: page table manipulation, interrupt
handling, and privileged register access.  The process can manage
its own virtual memory, handle its own page faults, and implement
custom exception handlers, all without modifying the host kernel
(beyond loading the Dune kernel module).

Dune's design is motivated by applications that benefit from
hardware privilege features:

\begin{itemize}[leftmargin=2em]
  \item \textbf{Sandboxing.}  A process can create a ring 3 sandbox
        within its own address space, with hardware-enforced isolation.
  \item \textbf{Garbage collection.}  A runtime can use dirty bits
        in page table entries to track memory writes, enabling
        write-barrier-free garbage collection.
  \item \textbf{Networking.}  A process can map NIC registers into
        its address space and handle interrupts directly, bypassing
        the kernel's network stack for zero-copy I/O.
\end{itemize}

\subsection{Relationship to UML's KVM backend}

Dune and \UML{}'s \KVMvTwo{} backend share a fundamental mechanism:
both use KVM to run code in a non-root mode VM with access to
privileged features.  The differences are in scope and purpose:

\begin{description}[style=nextline,leftmargin=2em]
  \item[Scope.]
    Dune gives privilege to a \emph{single process}.  \KVMvTwo{}
    runs an \emph{entire kernel} in the VM, with that kernel managing
    its own guest processes.

  \item[Kernel involvement.]
    In Dune, the application code running in ring 0 is arbitrary
    user-written code.  In \UML{}'s \KVMvTwo{}, the code running
    in the VM's ring 0 is the Linux kernel itself.

  \item[Process model.]
    Dune's VM contains a single application.  \UML{}'s VM contains
    a kernel that can create and schedule multiple guest processes,
    each running in the VM's ring 3.

  \item[Hardware requirements.]
    Both require \texttt{/dev/kvm} and VT-x.  \UML{} has fallback
    backends (seccomp, ptrace) that do not.
\end{description}

Dune can be viewed as a building block for what \UML{}'s \KVMvTwo{}
backend accomplishes at a larger scale.  Dune demonstrates that
KVM's non-root mode is useful for more than traditional VMs; \UML{}'s
\KVMvTwo{} backend demonstrates that it can host an entire kernel.


% ------------------------------------------------------------------
\section{Library OS: Graphene / Gramine}
\label{sec:related-libos}
% ------------------------------------------------------------------

\subsection{Architecture}

The Library OS concept~\cite{porter2011} proposes moving OS
functionality out of the kernel and into a library linked with the
application.  Instead of making system calls to a shared kernel,
the application calls library functions that implement OS
semantics.  The library OS runs on a minimal ``PAL'' (Platform
Adaptation Layer) that provides basic primitives: memory mapping,
thread creation, and I/O.

Graphene~\cite{tsai2014}, later renamed Gramine~\cite{gramine},
is the most prominent implementation of this concept.  Gramine
runs unmodified Linux binaries by intercepting their syscalls and
handling them within a library OS that runs in the same process.
Gramine's primary use case is running applications inside Intel
SGX enclaves, where the untrusted host OS cannot be relied upon
to correctly implement system calls.

\subsection{Comparison with UML}

\begin{description}[style=nextline,leftmargin=2em]
  \item[Shared goal.]
    Both Gramine and \UML{} aim to run Linux applications in
    environments where the host kernel is not directly serving
    the application's system calls.  Gramine interposes a library;
    \UML{} interposes an entire kernel.

  \item[Compatibility.]
    Gramine implements a subset of the Linux ABI (approximately
    150 syscalls).  Like gVisor, it faces the reimplementation
    problem: subtle differences from the kernel's behavior cause
    application failures.  \UML{} implements the complete Linux
    ABI because it is the complete Linux kernel.

  \item[Threat model.]
    Gramine's threat model assumes an untrusted host OS (SGX
    enclaves).  \UML{}'s threat model assumes a trusted host
    kernel.  These are essentially opposite assumptions.

  \item[Performance.]
    Gramine's in-process syscall handling is fast for
    compute-intensive workloads but incurs overhead for I/O
    (which must be forwarded through the PAL to the untrusted
    host).  \UML{}'s overhead is primarily in syscall interception
    (100--300\,ns per syscall with modern backends).

  \item[Use case.]
    Gramine is for running applications in hostile environments
    (SGX, cloud with untrusted host).  \UML{} is for running
    kernels in development/test environments.  The overlap is
    minimal.
\end{description}


% ------------------------------------------------------------------
\section{WSL1}
\label{sec:related-wsl1}
% ------------------------------------------------------------------

\subsection{Architecture}

Windows Subsystem for Linux version 1~\cite{wsl1}, shipped in
Windows 10 Anniversary Update (2016), implemented a Linux-compatible
kernel interface on Windows through syscall translation.  When a
Linux binary executed a syscall, the Windows kernel's \emph{lxcore}
driver translated it into equivalent Windows NT kernel calls.

WSL1 was a remarkable engineering achievement: it could run unmodified
Linux ELF binaries on a Windows kernel without a VM.  For many
workloads (development tools, scripting, basic server software),
it worked well.

\subsection{The reimplementation problem}

WSL1's approach is conceptually identical to gVisor's: reimplement
the Linux syscall interface on a different substrate.  And it
encountered the same fundamental problem: the long tail of
compatibility issues.  Syscalls that worked differently in subtle
ways (\texttt{fork} semantics, signal delivery timing,
\texttt{/proc} filesystem contents, inotify behavior) caused
failures in applications that depended on exact Linux semantics.

Microsoft ultimately abandoned the syscall-translation approach in
favor of WSL2, which runs a real Linux kernel in a Hyper-V virtual
machine.  This decision is a powerful validation of the principle
that \UML{} embodies: if you need Linux semantics, run the Linux
kernel.

\subsection{Lessons for UML}

WSL1's trajectory---from ambitious syscall translation to pragmatic
VM hosting---illustrates the fundamental difficulty of reimplementing
a kernel interface.  The Linux kernel has accumulated 35 years of
syscall semantics, much of it undocumented, some of it accidental,
all of it depended upon by real applications.  Any reimplementation
will eventually encounter the long tail of edge cases where it
diverges from the real kernel.

\UML{} avoids this problem entirely.  It is not a reimplementation;
it is the real kernel running on a different substrate.  The
compatibility question for \UML{} is not ``does it implement
\texttt{epoll} correctly?'' but ``does it intercept and forward
syscalls correctly?''---a much smaller and more tractable problem.


% ------------------------------------------------------------------
\section{Comprehensive Comparison}
\label{sec:related-table}
% ------------------------------------------------------------------

Table~\ref{tab:related-comparison} presents a comprehensive
comparison of the systems discussed in this chapter across seven
dimensions: approach, year, kernel compatibility, performance
overhead, hardware requirements, primary use case, and security
model.

\begin{table}[ht]
\centering
\caption{Comprehensive comparison of related systems}
\label{tab:related-comparison}
\small
\renewcommand{\arraystretch}{1.3}
\begin{tabularx}{\textwidth}{l l L{2.5cm} L{2.0cm} L{1.8cm} L{2.5cm}}
\toprule
\textbf{System} & \textbf{Year} & \textbf{Approach} & \textbf{Kernel compat.} & \textbf{HW req.} & \textbf{Primary use} \\
\midrule
\textbf{UML}
  & 2000
  & Kernel as process
  & Full (is the kernel)
  & None\textsuperscript{a}
  & Dev, debug, test \\
gVisor
  & 2018
  & Syscall reimpl.\ (Go)
  & Partial ($\sim$200 syscalls)
  & None\textsuperscript{a}
  & Container sandbox \\
Firecracker
  & 2018
  & MicroVM (KVM)
  & Full (separate kernel)
  & VT-x/AMD-V
  & Multi-tenant cloud \\
\textsc{LKL}
  & 2010
  & Kernel as library
  & Partial (no processes)
  & None
  & Fuzzing, research \\
QEMU/KVM
  & 2005/07
  & Full VM + HW virt
  & Full (any arch)
  & VT-x\textsuperscript{b}
  & General purpose \\
Containers
  & 2007+
  & Shared kernel
  & Full (same kernel)
  & None
  & App deployment \\
Unikernels
  & 2013
  & Single-image OS
  & Minimal
  & Hypervisor
  & Specialized services \\
Dune
  & 2012
  & User-level VT-x
  & N/A (not a kernel)
  & VT-x
  & Process sandboxing \\
Gramine
  & 2014
  & Library OS (SGX)
  & Partial ($\sim$150 syscalls)
  & SGX\textsuperscript{c}
  & Confidential computing \\
WSL1
  & 2016
  & Syscall translation
  & Partial (long tail)
  & None
  & Dev tools on Windows \\
\bottomrule
\end{tabularx}

\smallskip
\noindent
\textsuperscript{a}\,With seccomp/ptrace backend; KVM backend requires /dev/kvm. \quad
\textsuperscript{b}\,QEMU TCG mode does not require VT-x but is 10--100$\times$ slower. \quad
\textsuperscript{c}\,Gramine also runs without SGX in ``direct'' mode.
\end{table}

Table~\ref{tab:related-properties} compares the systems across
qualitative dimensions relevant to \UML{}'s target use cases.

\begin{table}[ht]
\centering
\caption{Qualitative property comparison}
\label{tab:related-properties}
\small
\renewcommand{\arraystretch}{1.3}
\begin{tabularx}{\textwidth}{l c c c c c}
\toprule
\textbf{System}
  & \textbf{Native gdb}
  & \textbf{Kernel isolation}
  & \textbf{Fast boot}
  & \textbf{Record/replay}
  & \textbf{Snapshot} \\
\midrule
\textbf{UML (v2)}
  & \checkmark & \checkmark & \checkmark & \checkmark & \checkmark \\
gVisor
  & ---        & \checkmark\textsuperscript{*} & \checkmark & --- & --- \\
Firecracker
  & ---        & \checkmark & \checkmark & --- & \checkmark \\
\textsc{LKL}
  & \checkmark & ---        & N/A        & ---        & --- \\
QEMU/KVM
  & ---        & \checkmark & ---        & ---        & \checkmark \\
Containers
  & ---        & ---        & \checkmark & ---        & --- \\
Unikernels
  & ---        & \checkmark & \checkmark & ---        & --- \\
Dune
  & ---        & ---        & N/A        & ---        & --- \\
Gramine
  & ---        & \checkmark\textsuperscript{*} & \checkmark & --- & --- \\
WSL1
  & ---        & \checkmark\textsuperscript{*} & \checkmark & --- & --- \\
\bottomrule
\end{tabularx}

\smallskip
\noindent
\textsuperscript{*}\,Isolation through reimplementation (not a real kernel), so
kernel-level bugs in the host are not exposed but compatibility gaps exist.\\
\checkmark\,= supported. \quad ---\,= not supported or not applicable.
\end{table}


% ------------------------------------------------------------------
\section{Architectural Taxonomy}
\label{sec:related-taxonomy}
% ------------------------------------------------------------------

The systems surveyed in this chapter can be organized along two
primary architectural axes: the \emph{execution substrate} (how
guest code is run) and the \emph{kernel model} (what provides
OS semantics).

\begin{figure}[ht]
\centering
\begin{tikzpicture}[
    sys/.style={
      draw=UMLGrid,
      fill=#1,
      minimum width=2.2cm,
      minimum height=0.8cm,
      rounded corners=3pt,
      font=\small\bfseries
    }
  ]

  % Axes
  \draw[-{Stealth[length=6pt]},thick,color=UMLMuted]
    (-1,-0.5) -- (-1,7) node[above,font=\small] {Kernel fidelity};
  \draw[-{Stealth[length=6pt]},thick,color=UMLMuted]
    (-1,-0.5) -- (11,-0.5) node[right,font=\small] {Isolation strength};

  % Systems positioned
  \node[sys=UMLRed!10]    (lkl)   at (0.5,0.8)  {LKL};
  \node[sys=UMLAmber!10]  (dune)  at (3.5,1.5)  {Dune};
  \node[sys=UMLAmber!10]  (gvis)  at (5.5,3.0)  {gVisor};
  \node[sys=UMLAmber!10]  (gram)  at (7.5,2.5)  {Gramine};
  \node[sys=UMLAmber!10]  (wsl)   at (3.0,3.5)  {WSL1};
  \node[sys=KernelGold!15](cont)  at (1.5,5.0)  {Containers};
  \node[sys=UMLBlue!10]   (fc)    at (8.5,5.5)  {Firecracker};
  \node[sys=UMLBlue!10]   (qemu)  at (9.5,6.3)  {QEMU/KVM};
  \node[sys=UMLTeal!10]   (unik)  at (6.5,1.0)  {Unikernels};
  \node[sys=UMLGreen!15]  (uml)   at (5.5,5.8)  {\UML{} v2};

  % Annotations
  \node[font=\scriptsize\itshape,color=UMLMuted] at (1.0,7.0)
    {Real Linux kernel};
  \node[font=\scriptsize\itshape,color=UMLMuted] at (1.0,0.0)
    {Reimplementation / none};
  \node[font=\scriptsize\itshape,color=UMLMuted] at (0.5,-1.0)
    {No boundary};
  \node[font=\scriptsize\itshape,color=UMLMuted] at (10.5,-1.0)
    {Hardware VM};

\end{tikzpicture}
\caption[Architectural taxonomy of isolation systems]{%
  Two-dimensional taxonomy of the systems discussed in this chapter.
  The vertical axis measures kernel fidelity (how closely the system
  matches real Linux kernel behavior).  The horizontal axis measures
  isolation strength (the degree of hardware/software boundary between
  the guest and host).  \UML{} v2 occupies a unique position: high
  kernel fidelity (it is the real kernel) with strong software-enforced
  isolation.
}
\label{fig:related-taxonomy}
\end{figure}

\subsection{Execution substrate}

\begin{description}[style=nextline,leftmargin=2em]
  \item[Hardware VM (QEMU/KVM, Firecracker).]
    Guest code runs on real CPU hardware in a hardware-isolated VM.
    Maximum isolation, maximum compatibility with real hardware, but
    requires VT-x/AMD-V and imposes VM exit overhead.

  \item[Process-level interposition (UML, gVisor).]
    Guest code runs as host processes.  Syscalls are intercepted
    via ptrace, seccomp, or KVM and handled by a supervisor
    (the UML kernel or gVisor's Sentry).  No hardware
    virtualization required (in non-KVM modes).

  \item[In-process / library (LKL, unikernels, Gramine).]
    Application and kernel code share an address space.  No
    interposition overhead, but no isolation between application
    and kernel.

  \item[Kernel extension (containers, Dune).]
    The host kernel is extended or configured to provide isolation
    (containers via namespaces/cgroups, Dune via VT-x non-root mode).
    Guest code runs directly on the host kernel with restricted
    visibility.
\end{description}

\subsection{Kernel model}

\begin{description}[style=nextline,leftmargin=2em]
  \item[Real Linux kernel (UML, QEMU/KVM, Firecracker, containers).]
    The real Linux kernel provides OS semantics.  In containers, it
    is the host kernel; in QEMU/KVM and Firecracker, it is a
    guest kernel; in \UML{}, it is the \UML{} kernel binary.

  \item[Reimplemented kernel (gVisor, WSL1, Gramine).]
    A separate codebase reimplements (a subset of) the Linux
    syscall interface.  This avoids exposing the real kernel's
    attack surface but introduces compatibility risk.

  \item[Minimal / absent kernel (LKL, unikernels).]
    The kernel is either compiled as a library (LKL) or reduced
    to the minimum needed for a specific application (unikernels).
    Kernel semantics are partial or specialized.
\end{description}

\UML{} is unique in combining the real Linux kernel (maximum
compatibility and fidelity) with process-level execution (host-tool
accessibility and minimal hardware requirements).  This combination
is what makes it irreplaceable for the use cases described in
Chapter~\ref{ch:applications}.


% ------------------------------------------------------------------
\section{Summary}
\label{sec:related-summary}
% ------------------------------------------------------------------

The related work surveyed in this chapter demonstrates that \UML{}
occupies a distinct and valuable position in the isolation/virtualization
landscape.  The key insight is that most related systems optimize for
\emph{either} production deployment (Firecracker, gVisor, containers,
unikernels) \emph{or} development speed (LKL, Dune), but not both
kernel fidelity and developer experience.

\UML{}'s contribution is precisely this combination:

\begin{enumerate}[leftmargin=2em]
  \item \textbf{Full kernel fidelity.}  Unlike gVisor, WSL1, and
        Gramine, \UML{} does not reimplement the kernel.  A bug
        found in \UML{} is a bug in Linux.

  \item \textbf{Full user/kernel separation.}  Unlike \textsc{LKL}
        and unikernels, \UML{} maintains the process model, scheduling,
        and memory protection.

  \item \textbf{No hardware requirements.}  Unlike QEMU/KVM,
        Firecracker, and Dune, \UML{} (with seccomp) runs wherever
        a Linux process can run.

  \item \textbf{Native host-tool access.}  Unlike every VM-based
        system, \UML{} exposes the kernel to \texttt{gdb},
        \texttt{strace}, \texttt{perf}, and every other host tool
        as a first-class debugging target.

  \item \textbf{Kernel isolation.}  Unlike containers, \UML{} provides
        per-instance kernel isolation, making it safe for security
        research and untrusted workloads.
\end{enumerate}

The v2 redesign strengthens \UML{}'s position by adding capabilities
(snapshot/restore, record/replay, SMP) that were previously available
only in heavier systems (QEMU, Firecracker) or not available at all
(kernel-level record/replay).  Chapter~\ref{ch:v2redesign} describes
these contributions in detail.
